% C2_brief.tex -- ONE source for the English and Romanian case brief of Case 2 (Lecture 2).
\input{numbers_\lang}
\renewcommand{\emqacase}{\T{Research case 2 (Lecture 2)}{Studiul de caz 2 (cursul 2)}}
\begin{document}
\emqatitle{\T{Case 2: Day-ahead forecasting of national electricity consumption}%
                {Studiul de caz 2: Prognoza consumului național de energie electrică pentru ziua următoare}}%
          {\T{Based on}{Studiu bazat pe lucrarea} Proșcanu, Proșcanu, Pele \& Costea (2025), \emph{Watt's next? Benchmarking time series models on Romania's national electricity consumption},
           Romanian Journal of Economic Forecasting 28(3), 146--163, \href{https://EconPapers.repec.org/RePEc:rjr:romjef:v::y:2025:i:3:p:146-163}{RePEc}}

\begin{scenario}
\T{Every morning, a Romanian supplier or the balancing team of a producer must notify the transmission system operator (TSO), Transelectrica, of the volume it will consume or deliver in each hour of the next day.
Every MWh of error becomes an \emph{imbalance}, settled at a price that is usually worse than the day-ahead price.
The head of trading must decide: \textbf{which forecasting model do we run in production --- a classical seasonal model, a machine-learning model, or a pretrained ``foundation'' model --- and is the gain worth the cost?}}%
{În fiecare dimineață, un furnizor din România sau echipa de echilibrare a unui producător trebuie să transmită operatorului de transport și de sistem (OTS), Transelectrica, notificarea privind volumul pe care îl va consuma sau livra în fiecare oră din ziua următoare.
Fiecare MWh de eroare se transformă într-un \emph{dezechilibru}, decontat de regulă la un preț mai dezavantajos decît cel de pe piața pentru ziua următoare (PZU).
Directorul de tranzacționare trebuie să decidă: \textbf{ce model de prognoză folosim în activitatea curentă --- un model sezonier clasic, un model de învățare automată sau un model fundamental preantrenat (\emph{foundation model}) --- și justifică îmbunătățirea acurateței costul implementării?}}
\end{scenario}

\section{\T{What the paper did}{Lucrarea de referință}}
\textbf{\T{Question.}{Întrebarea de cercetare.}} \T{The paper benchmarks univariate forecasting models on Romania's national electricity consumption, (1) on predictive accuracy (MSE, MAE, $R^2$) and (2) on the effect of model architecture (p.~147).}%
{Lucrarea compară modele de prognoză univariate aplicate consumului național de energie electrică al României, (1) din punctul de vedere al acurateței (MSE, MAE, $R^2$) și (2) din punctul de vedere al efectului pe care îl are arhitectura modelului (p.~147).}

\textbf{\T{Data.}{Datele.}} \T{Hourly national consumption, 1 January 2022 -- 31 December 2024, 26,304 observations in MWh (sistemulenergetic.ro); mean 6,230.59, minimum 2,871.75, maximum 9,210.75 (Table~1, p.~149).
Lowest demand at 2--5 a.m., peak around 7 p.m.; winter months highest (p.~150).}%
{Consumul național orar, 1 ianuarie 2022 -- 31 decembrie 2024, 26.304 observații, în MWh (sistemulenergetic.ro); media 6.230,59, minimul 2.871,75, maximul 9.210,75 (Tabelul~1, p.~149).
Consumul este minim între orele 2 și 5 și atinge vîrful în jurul orei 19; lunile de iarnă au cel mai mare consum (p.~150).}

\textbf{\T{Method.}{Metoda.}} \T{Nine models: SARIMA (Seasonal AutoRegressive Integrated Moving Average), the neural networks N-BEATS, NHITS, LSTM (Long Short-Term Memory), Mamba, Titans and RetNet,
and two pretrained foundation models used zero-shot, Chronos-T5 and TimesFM. Look-back 168 hours, horizon 24 hours, chronological 80/20 split (Table~3, p.~157).}%
{Nouă modele: SARIMA (model autoregresiv integrat cu medie mobilă, cu componentă sezonieră), rețelele neuronale N-BEATS, NHITS, LSTM (\emph{Long Short-Term Memory}, rețea recurentă cu memorie de lungă durată), Mamba, Titans și RetNet,
precum și două modele fundamentale preantrenate, folosite fără reantrenare pe datele românești (\emph{zero-shot}): Chronos-T5 și TimesFM. Fereastra de intrare are 168 de ore, orizontul de prognoză 24 de ore, iar eșantionul este împărțit cronologic în proporția 80/20 (Tabelul~3, p.~157).}

\textbf{\T{Main findings of the paper.}{Principalele rezultate ale lucrării.}}
\begin{itemize}
\item \T{``The N-BEATS model has been the most successful'', with MAPE 2.324\,\% and ``the highest coefficient of determination ($R^2$ = 0.956)'' (p.~158; Table~4, p.~157).
N-BEATS ``lowers MAPE from 3.63\,\% (SARIMA) to 2.32\,\%'' (p.~159).}%
{Cel mai performant model este N-BEATS, cu MAPE 2,324\,\% și cel mai mare coeficient de determinare ($R^2 = 0{,}956$) (p.~158; Tabelul~4, p.~157).
Față de SARIMA, N-BEATS reduce MAPE de la 3,63\,\% la 2,32\,\% (p.~159).}
\item \T{RetNet and Titans also perform very well (p.~158--159).}%
{RetNet și Titans au, de asemenea, performanțe foarte bune (p.~158--159).}
\item \T{``TimesFM consistently underperforms'' (MAPE 5.895\,\%, $R^2$ = 0.673); Chronos-T5 is the fastest (7 seconds) with MAPE 3.325\,\% (p.~159; Table~4).}%
{TimesFM are în mod constant rezultate mai slabe (MAPE 5,895\,\%, $R^2 = 0{,}673$); Chronos-T5 este cel mai rapid (7 secunde), cu MAPE 3,325\,\% (p.~159; Tabelul~4).}
\item \T{SARIMA ``remains a viable and efficient option'' when resources are constrained; neural models suit capacity planning and renewable integration (p.~160).}%
{SARIMA rămîne o opțiune viabilă și eficientă atunci cînd resursele de calcul sînt limitate; modelele neuronale sînt potrivite pentru planificarea capacităților și integrarea surselor regenerabile (p.~160).}
\end{itemize}

\section{\T{The method, step by step}{Metoda, pas cu pas}}
\begin{intuition}
\T{Consumption repeats itself: every day has the same shape, every week has quiet weekends, every year has cold winters. A forecasting model is a machine for copying the right past pattern and adjusting it.
Simple models copy one pattern (``two days ago'', ``last week''); better models learn which past hours matter and how much.}%
{Consumul are un caracter repetitiv: profilul orar se repetă de la o zi la alta, sfîrșitul de săptămînă aduce un consum mai redus, iar iernile sînt reci în fiecare an. Un model de prognoză preia tiparul potrivit din trecut și îl ajustează.
Modelele simple preiau un singur tipar („acum două zile”, „săptămîna trecută”); modelele mai bune învață care ore din trecut sînt relevante și cu ce pondere.}
\end{intuition}
\T{Let $L_{d,h}$ be the load (MW) in hour $h$ of day $d$. At gate closure, 12:00 CET on day $d-1$, we forecast the 24 values $\hat L_{d,0},\dots,\hat L_{d,23}$. At that time load is metered only up to the end of day $d-2$, so:
\begin{itemize}\setlength{\itemsep}{0pt}
\item every model uses load data up to day $d-2$ only (lags of at least 48 hours, plus $d-7$);
\item no value of day $d-1$ is used, not even its daily mean.
\end{itemize}
The \emph{naive} forecasts are $\hat L_{d,h}=L_{d-2,h}$ (two days earlier) and $\hat L_{d,h}=L_{d-7,h}$ (last week).
SARIMA$(1,0,1)(1,1,1)_{24}$ models the 24-hour difference $L_t-L_{t-24}$ as depending on its own past and on past errors.
The \emph{linear regression} and \emph{gradient boosting} (an ensemble of small decision trees, each fitted to the errors of the previous ones) use lagged load (48, 72, 168, 336 hours), the mean of day $d-2$, hour, weekday, month and holidays.
The \emph{neural network} maps the 168 hours up to the end of $d-2$ directly to the 24 hours of day $d$.
Chronos-Bolt is a model pretrained on many public time series; it reads the 512 hours up to the end of $d-2$ and outputs forecast quantiles without being trained on our data (zero-shot).
Accuracy is measured on all $N$ test hours with}%
{Fie $L_{d,h}$ consumul (MW) în ora $h$ a zilei $d$. La închiderea porții, ora 12:00 CET în ziua $d-1$, prognozăm cele 24 de valori $\hat L_{d,0},\dots,\hat L_{d,23}$. În acel moment, consumul este măsurat numai pînă la sfîrșitul zilei $d-2$, de aceea:
\begin{itemize}\setlength{\itemsep}{0pt}
\item fiecare model folosește date de consum numai pînă în ziua $d-2$ (decalaje de cel puțin 48 de ore, plus $d-7$);
\item nu se folosește nicio valoare din ziua $d-1$, nici măcar media ei zilnică.
\end{itemize}
Prognozele \emph{naive} sînt $\hat L_{d,h}=L_{d-2,h}$ (cu două zile înainte) și $\hat L_{d,h}=L_{d-7,h}$ (aceeași zi din săptămîna precedentă).
Modelul SARIMA$(1,0,1)(1,1,1)_{24}$ descrie diferența pe 24 de ore, $L_t-L_{t-24}$, în funcție de valorile și de erorile sale anterioare.
\emph{Regresia liniară} și modelul \emph{gradient boosting} (un ansamblu de arbori de decizie mici, fiecare estimat pe erorile celor anteriori) folosesc consumul decalat (cu 48, 72, 168 și 336 de ore), media zilei $d-2$, ora, ziua săptămînii, luna și zilele de sărbătoare.
\emph{Rețeaua neuronală} transformă direct cele 168 de ore de pînă la sfîrșitul zilei $d-2$ în prognoza pentru cele 24 de ore ale zilei $d$.
Chronos-Bolt este un model preantrenat pe un număr mare de serii de timp publice; preia cele 512 ore de pînă la sfîrșitul zilei $d-2$ și furnizează cuantilele prognozei fără a fi antrenat pe datele noastre (\emph{zero-shot}).
Acuratețea se măsoară pe toate cele $N$ ore din perioada de test, cu indicatorii}
\[
\text{MAE}=\frac1N\sum_{t}|L_t-\hat L_t|,\qquad \text{MAPE}=\frac{100}{N}\sum_t\left|\frac{L_t-\hat L_t}{L_t}\right|,\qquad R^2=1-\frac{\sum_t(L_t-\hat L_t)^2}{\sum_t(L_t-\bar L)^2}.
\]
\begin{decision}
\T{The model choice is a money question: annual imbalance volume $\approx$ MAE $\times$ 8,760 hours $\times$ your market share; multiply by the imbalance cost per MWh. A model is worth deploying if this saving exceeds its running and maintenance cost.}%
{Alegerea modelului are o miză financiară: volumul anual al dezechilibrelor $\approx$ MAE $\times$ 8.760 de ore $\times$ cota dumneavoastră de piață; înmulțit cu costul dezechilibrului pe MWh, acesta dă costul anual. Un model merită introdus în activitatea curentă dacă economia astfel obținută depășește costurile de operare și de întreținere.}
\end{decision}

\section{\T{Replicate it yourself}{Replicare pe datele cursului}}
\T{Notebook \texttt{C2\_notebook\_EN.ipynb}: hourly load from ENTSO-E via energy-charts.info (\texttt{generation\_ro\_mw\_hourly\_utc.csv}), the paper's sample and an 80/20 split; the test period starts on \NtestStart\ (\NnDays\ forecast days).
Our series reproduces the paper's Table~1: \NnHours\ hours, mean \NmeanLoad\ MW, minimum \NminLoad, maximum \NmaxLoad\ --- the same Transelectrica metering, published twice.}%
{Notebook-ul \texttt{C2\_notebook\_RO.ipynb} folosește consumul orar publicat de ENTSO-E, preluat de pe energy-charts.info (\texttt{generation\_ro\_mw\_hourly\_utc.csv}), același eșantion ca lucrarea și aceeași împărțire 80/20; perioada de test începe la \NtestStart\ (\NnDays\ zile de prognoză).
Seria noastră reproduce Tabelul~1 din lucrare: \NnHours\ ore, media \NmeanLoad\ MW, minimul \NminLoad, maximul \NmaxLoad\ --- sînt aceleași măsurători Transelectrica, publicate de două surse.}

\begin{figure}[H]
\begin{minipage}[t]{0.55\linewidth}\centering
\includegraphics[width=\linewidth]{f1_load_series}
\caption{\T{Daily mean load and the test period (left); mean hourly profile by month (right).}{Consumul mediu zilnic și perioada de test (stînga); profilul orar mediu, pe luni (dreapta).}}
\src{energy-charts.info (ENTSO-E); 1.1.2022--31.12.2024.}
\end{minipage}\hfill
\begin{minipage}[t]{0.43\linewidth}\centering
\includegraphics[width=\linewidth]{f2_mape_by_model}
\caption{\T{Test MAPE by model.}{MAPE pe perioada de test, pe modele.}}
\src{\T{our forecasts; test period \NtestStart--31.12.2024.}{prognozele noastre; perioada de test \NtestStart--31.12.2024.}}
\end{minipage}
\end{figure}

\begin{table}[H]\centering\small
\caption{\T{Day-ahead forecast accuracy on the test period, \NnDays\ days $\times$ 24 hours. Best value in bold (Bias closest to zero).}%
{Acuratețea prognozelor pentru ziua următoare în perioada de test, \NnDays\ zile $\times$ 24 de ore. Cea mai bună valoare este marcată cu aldine (pentru abatere, valoarea cea mai apropiată de zero).}}
\tabin{accuracy}
\src{\T{our computation; Bias = mean of actual minus forecast.}{calcule proprii; abaterea = media diferenței dintre valoarea efectivă și prognoză.}}
\end{table}

\begin{figure}[H]
\begin{minipage}[t]{0.55\linewidth}\centering
\includegraphics[width=\linewidth]{f3_forecast_week}
\caption{\T{One week of forecasts in November 2024.}{O săptămînă de prognoze în noiembrie 2024.}}
\src{\T{our forecasts; energy-charts.info.}{prognozele noastre; energy-charts.info.}}
\end{minipage}\hfill
\begin{minipage}[t]{0.43\linewidth}\centering
\includegraphics[width=\linewidth]{f4_error_by_hour}
\caption{\T{MAE by hour of the forecast day.}{MAE pe orele zilei prognozate.}}
\src{\T{our forecasts; test period.}{prognozele noastre; perioada de test.}}
\end{minipage}
\end{figure}

\T{\textbf{Our result.} The best model on our test period is \NbestName\ with MAPE \NbestMAPE\,\% (MAE \NbestMAE\ MW, $R^2$ = \NbestRtwo), followed closely by Chronos-Bolt used zero-shot (\NchronosMAPE\,\%).
The naive ``two days earlier'' rule has \NnaiveMAPE\,\%, ``last week'' \NweekMAPE\,\%; our SARIMA reaches \NsarimaMAPE\,\%, with parameters estimated on 8 weeks and held fixed over the test, so it does not adapt to holidays and the summer.
As in the paper, a pretrained model used zero-shot is competitive, and the ranking depends on the metric. The 10\,\%--90\,\% interval of Chronos-Bolt covered \Ncover\,\% of the hours (target 80\,\%).
Our numbers reflect our test window, forecast origin and model versions; the ranking is ex post (the ``best'' model is identified on the test period itself), so the saving below is an upper estimate.
\textbf{Decision number:} for a supplier with 10\,\% of national consumption and an illustrative imbalance cost of 50 EUR/MWh, moving from the naive rule to the best model cuts the annual imbalance bill from \NcostNaive\ to \NcostBest\ million EUR (saving \Nsaving\ million EUR).}%
{\textbf{Rezultatul replicării.} În perioada noastră de test, cel mai bun model este \NbestName, cu MAPE \NbestMAPE\,\% (MAE \NbestMAE\ MW, $R^2 = \NbestRtwo$), urmat îndeaproape de Chronos-Bolt, folosit fără reantrenare (\NchronosMAPE\,\%).
Regula naivă „acum două zile” are MAPE \NnaiveMAPE\,\%, iar regula „săptămîna trecută” \NweekMAPE\,\%; modelul SARIMA ajunge la \NsarimaMAPE\,\%, deoarece parametrii săi sînt estimați pe 8 săptămîni și rămîn neschimbați pe toată perioada de test, astfel că nu se adaptează la sărbători și la sezonul de vară.
Ca și în lucrare, un model preantrenat folosit fără reantrenare este competitiv, iar ierarhia modelelor depinde de indicatorul de eroare ales. Intervalul dintre cuantilele de 10\,\% și 90\,\% ale modelului Chronos-Bolt a cuprins consumul efectiv în \Ncover\,\% din ore (nivelul-țintă este 80\,\%).
Rezultatele depind de perioada de test, de momentul prognozei și de versiunile modelelor folosite; ierarhia este stabilită ex post („cel mai bun” model este identificat chiar pe perioada de test), de aceea economia de mai jos este o estimare superioară.
\textbf{Cifra relevantă pentru decizie:} pentru un furnizor cu 10\,\% din consumul național și un cost ilustrativ al dezechilibrului de 50 EUR/MWh, trecerea de la regula naivă la cel mai bun model reduce costul anual al dezechilibrelor de la \NcostNaive\ la \NcostBest\ milioane EUR (o economie de \Nsaving\ milioane EUR).}

\section{\T{Discussion questions for managers}{Întrebări pentru discuție}}
\begin{enumerate}
\item \T{Using the paper's benchmarking framework, what evidence would you need before choosing a model for production: seasonal performance, imbalance cost, runtime or maintenance?}%
{Pornind de la cadrul de comparație din lucrare, de ce informații ați avea nevoie înainte de a alege modelul care va fi folosit în activitatea curentă: performanța pe sezoane, costul dezechilibrelor, timpul de calcul sau efortul de întreținere?}
\item \T{A pretrained model needs no training on your data. What are the risks of relying on a model you cannot inspect or retrain?}%
{Un model preantrenat nu trebuie antrenat pe datele dumneavoastră. Ce riscuri presupune să vă bazați pe un model pe care nu îl puteți analiza în detaliu sau reantrena?}
\item \T{Which hours of the day carry the largest errors (Figure~4), and are they the expensive hours for your portfolio?}%
{În ce ore ale zilei apar cele mai mari erori (Figura~4) și sînt acestea orele cele mai costisitoare pentru portofoliul dumneavoastră?}
\item \T{National load is smooth; your own portfolio (a few hundred clients, rooftop solar) is not. How would you expect the ranking of models to change?}%
{Consumul național are o evoluție netedă; portofoliul dumneavoastră (cîteva sute de clienți, panouri fotovoltaice pe acoperișuri) nu are. Cum credeți că s-ar modifica ierarhia modelelor?}
\item \T{Is a MAPE improvement of one percentage point worth a data-science team? Use the decision number above and your own share and imbalance cost.}%
{Justifică o reducere a MAPE cu un punct procentual costul unei echipe de analiză a datelor? Folosiți cifra relevantă pentru decizie de mai sus, cu propria cotă de piață și propriul cost al dezechilibrului.}
\end{enumerate}

\section{\T{Limitations of the classroom replication}{Limitele replicării}}
\T{(i) Our models are univariate: no weather or economic drivers.
(ii) One test window (late May -- December 2024); rankings can change with the season.
(iii) National load hides the behind-the-meter solar boom; a supplier's own portfolio is noisier.
(iv) The imbalance cost of 50 EUR/MWh is an illustrative assumption, not a market statistic.
(v) Our neural network and Chronos version are compact implementations chosen to run in class.}%
{(i) Modelele noastre sînt univariate: nu includ variabile meteorologice sau economice.
(ii) Folosim o singură perioadă de test (sfîrșitul lunii mai -- decembrie 2024); ierarhia modelelor se poate schimba de la un sezon la altul.
(iii) Consumul național nu reflectă direct creșterea rapidă a producției fotovoltaice a prosumatorilor; portofoliul unui furnizor are o variabilitate mai mare.
(iv) Costul dezechilibrului de 50 EUR/MWh este o ipoteză ilustrativă, nu o statistică de piață.
(v) Rețeaua neuronală și versiunea Chronos folosite sînt implementări compacte, alese pentru a putea rula la curs.}

\newpage
\section*{\T{Summary of Case 2}{Studiul de caz 2 pe scurt}}
\glance{\T{Which model should forecast tomorrow's hourly consumption, and is the gain worth it?}{Ce model ar trebui să prognozeze consumul orar de mîine și justifică îmbunătățirea costul?}}%
{\T{Hourly national load 2022--2024 (\NnHours\ h), same series as the paper; last 20\,\% for testing.}{Consumul național orar 2022--2024 (\NnHours\ h), aceeași serie ca în lucrare; ultimele 20\,\% din date pentru test.}}%
{\T{Naive, SARIMA, regression, gradient boosting, neural net, Chronos zero-shot; MAE, MAPE, $R^2$.}{Naiv, SARIMA, regresie, gradient boosting, rețea neuronală, Chronos zero-shot; MAE, MAPE, $R^2$.}}%
{\T{Best: \NbestName, MAPE \NbestMAPE\,\% vs \NnaiveMAPE\,\% naive. Paper: N-BEATS 2.324\,\% vs SARIMA 3.63\,\%.}{Cel mai bun: \NbestName, MAPE \NbestMAPE\,\% față de \NnaiveMAPE\,\% pentru regula naivă. Lucrarea: N-BEATS 2,324\,\% față de SARIMA 3,63\,\%.}}%
{\T{Deploy a modern model if the imbalance saving ($\approx$\Nsaving\ mil.\ EUR/yr for a 10\,\% supplier) exceeds its cost.}{Introduceți un model modern dacă economia la costul dezechilibrelor ($\approx$\Nsaving\ mil.\ EUR/an pentru un furnizor cu 10\,\%) depășește costul acestuia.}}
\begin{center}
\includegraphics[width=0.62\linewidth]{fig_headline_\lang}
\src{\T{our forecasts on energy-charts.info (ENTSO-E) load; test period \NtestStart--31.12.2024.}{prognozele noastre pe datele de consum de pe energy-charts.info (ENTSO-E); perioada de test \NtestStart--31.12.2024.}}
\end{center}
\begin{minipage}[t]{0.32\linewidth}\centering{\color{MainBlue}\Huge\bfseries \NbestMAPE\,\%}\\ \small \T{best MAPE on our test}{cel mai bun MAPE în perioada de test}\end{minipage}\hfill
\begin{minipage}[t]{0.32\linewidth}\centering{\color{MainBlue}\Huge\bfseries \NnaiveMAPE\,\%}\\ \small \T{naive ``two days earlier'' rule}{regula naivă „acum două zile”}\end{minipage}\hfill
\begin{minipage}[t]{0.32\linewidth}\centering{\color{IDAred}\Huge\bfseries \Nsaving}\\ \small \T{mil.\ EUR/yr imbalance saving, 10\,\% supplier}{mil.\ EUR/an economie la costul dezechilibrelor, furnizor cu 10\,\%}\end{minipage}
\end{document}
